\documentclass[11pt,a4paper]{report}

\usepackage[margin=2.4cm]{geometry}
\usepackage[T1]{fontenc}
\usepackage[utf8]{inputenc}
\usepackage[dutch]{babel}
\usepackage{lmodern}
\usepackage{microtype}
\usepackage{hyperref}
\usepackage{booktabs}
\usepackage{longtable}
\usepackage{tabularx}
\usepackage{enumitem}
\usepackage{listings}
\usepackage{xcolor}
\usepackage{graphicx}
\usepackage{fancyhdr}
\usepackage{titlesec}
\usepackage{tocloft}
\usepackage{csquotes}

\hypersetup{
  colorlinks=true,
  linkcolor=black,
  urlcolor=blue!50!black,
  citecolor=black,
  pdftitle={AFC Ticketing --- Volledige technische documentatie},
  pdfauthor={AFC Gent Ticketing}
}

\definecolor{codebg}{RGB}{246,247,249}
\definecolor{codeframe}{RGB}{216,222,232}
\definecolor{brandnavy}{HTML}{081C3C}
\definecolor{brandorange}{HTML}{EC6525}

\lstset{
  basicstyle=\ttfamily\small,
  backgroundcolor=\color{codebg},
  frame=single,
  rulecolor=\color{codeframe},
  breaklines=true,
  columns=fullflexible,
  keepspaces=true,
  showstringspaces=false,
  tabsize=2,
  literate=
    {á}{{\'a}}1 {é}{{\'e}}1 {í}{{\'i}}1 {ó}{{\'o}}1 {ú}{{\'u}}1
    {Á}{{\'A}}1 {É}{{\'E}}1 {Í}{{\'I}}1 {Ó}{{\'O}}1 {Ú}{{\'U}}1
    {à}{{\`a}}1 {è}{{\`e}}1 {ì}{{\`i}}1 {ò}{{\`o}}1 {ù}{{\`u}}1
    {ä}{{\"a}}1 {ë}{{\"e}}1 {ï}{{\"i}}1 {ö}{{\"o}}1 {ü}{{\"u}}1
    {Ä}{{\"A}}1 {Ë}{{\"E}}1 {Ï}{{\"I}}1 {Ö}{{\"O}}1 {Ü}{{\"U}}1
    {â}{{\^a}}1 {ê}{{\^e}}1 {î}{{\^i}}1 {ô}{{\^o}}1 {û}{{\^u}}1
    {ç}{{\c c}}1 {Ç}{{\c C}}1
    {ñ}{{\~n}}1 {Ñ}{{\~N}}1
    {ß}{{\ss}}1
}

\pagestyle{fancy}
\fancyhf{}
\fancyhead[L]{\small AFC Ticketing --- technische documentatie}
\fancyhead[R]{\small\thepage}
\renewcommand{\headrulewidth}{0.4pt}

\title{\textbf{AFC Ticketing}\\[0.4em]
\large Volledige technische documentatie, architectuur en reconstructiegids}
\author{Gebaseerd op de codebase \texttt{ticketing} (Next.js 16 / React 19 / SQLite)}
\date{\today}

\begin{document}
\maketitle

\begin{abstract}
Dit document beschrijft de AFC Gent ticketing-applicatie tot op reconstructieniveau.
Het legt uit \emph{waarom} elke architectuurkeuze gemaakt is, \emph{wat} ieder bestand doet,
\emph{hoe} de datalaag, authenticatie, tickets, API's en gebruikersflows werken, en
\emph{hoe} je de tool vanaf nul opnieuw zou bouwen.
De applicatie is een monolithische Next.js~16 App Router-app met ingebouwde API-routes,
een lokale SQLite-database via \texttt{better-sqlite3}, optionele SMTP-mail, en HMAC-ondertekende
QR-tickets die ook op ephemeral serverless-schijven blijven werken.
\end{abstract}

\tableofcontents
\newpage

%======================================================================
\chapter{Doel, scope en ontwerpprincipes}
%======================================================================

\section{Wat de tool doet}
De applicatie ondersteunt meerdere events binnen één installatie:
\begin{itemize}[leftmargin=*]
  \item Publieke homepage met openstaande events (optioneel met foto).
  \item Inschrijfformulier per event (naam, e-mail, telefoon, voedselvoorkeur, optionele extra info, verplicht CV als PDF).
  \item Ticketpagina met QR-code na inschrijving; optionele bevestigingsmail.
  \item Adminportaal: login, dashboard, eventbeheer (CRUD/archief/foto/uitgebreide beschrijving),
        deelnemersbeheer (toevoegen/verwijderen/check-in), QR-scanner, CSV-export.
  \item Archiefpagina voor afgelopen events.
\end{itemize}

\section{Niet-functionele eisen die de architectuur sturen}
\begin{description}[style=nextline,leftmargin=0pt]
  \item[Eén process / één service]
    Geen aparte backend-API-server. Alles leeft in Next.js (\texttt{app/} voor UI, \texttt{app/api/} voor HTTP).
    Dit vermindert deployment-complexiteit voor een clubtool.
  \item[Locale persistentie]
    SQLite in \texttt{data/tickets.db} + uploads in \texttt{data/uploads/}. Geen Postgres-vereiste.
    Op Vercel is de schijf read-only buiten \texttt{/tmp}; daarom schrijft de runtime daar naartoe.
  \item[Serverless-robust tickets]
    Op Vercel heeft elke serverless-instance een eigen lege \texttt{/tmp}-SQLite.
    Daarom zijn tickets HMAC-ondertekend: de payload zelf bewijst geldigheid, zelfs als de lokale DB de rij mist.
  \item[Admin is ``good enough'', niet enterprise-IAM]
    Eén gedeeld wachtwoord + cookie-sessie. Geschikt voor een organisatiesleutel, niet voor multi-user RBAC.
  \item[Nederlandstalige UX]
    Copy, foutmeldingen, CSV-headers en datums (\texttt{nl-BE}, tijdzone \texttt{Europe/Brussels}).
\end{description}

\section{Waarom TypeScript}
\begin{itemize}[leftmargin=*]
  \item \textbf{Contracten over grenzen heen:} FormData $\to$ Zod $\to$ DB-row $\to$ camelCase domain model.
    TypeScript vangt mismatch tussen \texttt{RegistrationRow} (snake\_case SQL) en \texttt{Registration} (app-model).
  \item \textbf{Next.js App Router} is TypeScript-first (\texttt{params: Promise<...>}, route handlers, server/client boundary).
  \item \textbf{Native modules} zoals \texttt{better-sqlite3} hebben \texttt{@types}; compile-time checks vermijden runtime ``undefined column''.
  \item \textbf{Refactorveiligheid} bij multi-event migratie: slug, \texttt{event\_id}, signed claims, en admin FormData-velden blijven synchroon.
\end{itemize}
JavaScript zou werken, maar de kost van types is laag t.o.v.\ de kost van stille datafouten bij check-in en tickets.

\section{Waarom Next.js 16 (App Router) + React 19}
\begin{itemize}[leftmargin=*]
  \item Server Components voor data-reads dicht bij SQLite (homepage, adminlijsten) zonder aparte JSON-BFF.
  \item Client Components alleen waar nodig (formulieren, scanner, interactieve tabellen).
  \item Route Handlers als API zonder Express.
  \item \texttt{proxy.ts} (Next 16-conventie) beschermt \texttt{/admin/*} vóór page-render.
\end{itemize}
\textbf{Belangrijk:} deze Next-versie wijkt af van oudere docs. Lees altijd
\texttt{node\_modules/next/dist/docs/} in de repo vóór framework-wijzigingen (\texttt{AGENTS.md}).

\section{Waarom SQLite + better-sqlite3}
\begin{itemize}[leftmargin=*]
  \item Zero-ops voor een club: geen managed DB, backups = map \texttt{data/} kopiëren.
  \item \texttt{better-sqlite3} is synchroon en snel in Node --- ideaal in route handlers zonder connection pool.
  \item Nadelen bewust aanvaard: slechte fit voor multi-region serverless writes; Vercel-demo's zijn ephemeral.
    Productie op een VPS met persistente schijf is het bedoelde pad.
\end{itemize}

\section{Kernstack}
\begin{tabularx}{\textwidth}{@{}lX@{}}
\toprule
\textbf{Laag} & \textbf{Keuze} \\
\midrule
Runtime & Node 22, Next.js 16.3.1 (Turbopack in dev) \\
UI & React 19.2, Tailwind CSS 4, Google font Josefin Sans \\
Validatie & Zod 4 \\
DB & better-sqlite3 13, bestand \texttt{data/tickets.db} \\
QR & \texttt{qrcode} (server), \texttt{@yudiel/react-qr-scanner} (admin camera) \\
Mail & nodemailer (optioneel SMTP) \\
Auth & SHA-256 sessiecookie \texttt{gate\_session} \\
Tickets & HMAC-SHA256 tokenformaat \texttt{t1.<body>.<sig>} \\
\bottomrule
\end{tabularx}

%======================================================================
\chapter{Repositorystructuur --- elk bestand uitgelegd}
%======================================================================

\section{Overzicht boom}
\begin{lstlisting}
ticketing/
  app/                 # App Router: pages + API
  components/          # React UI (server/client)
  lib/                 # Domeinlogica, DB, auth, mail, tokens
  public/              # Statische assets (logo)
  data/                # Runtime: tickets.db + uploads/ (niet ``broncode'')
  proxy.ts             # Admin route guard
  next.config.ts
  package.json
  tsconfig.json
  AGENTS.md / README.md
\end{lstlisting}

\section{Rootconfiguratie}

\subsection{\texttt{package.json}}
Definieert scripts \texttt{dev}, \texttt{build}, \texttt{start}, \texttt{lint} en alle dependencies.
Er is \textbf{geen} testscript: kwaliteit wordt manueel/via lint+build afgedwongen.

\subsection{\texttt{tsconfig.json}}
Strict TypeScript, path alias \texttt{@/*} $\to$ projectroot. Daardoor imports als
\texttt{@/lib/db} in plaats van relatieve \texttt{../../}.

\subsection{\texttt{next.config.ts}}
Markeert native/CJS-modules (\texttt{better-sqlite3}, evt.\ nodemailer) correct voor bundling,
zodat server routes in de Node-runtime blijven.

\subsection{\texttt{postcss.config.mjs} + Tailwind 4}
PostCSS laadt \texttt{@tailwindcss/postcss}. Styling zit grotendeels in utility-classes plus
CSS-variabelen in \texttt{app/globals.css} (\texttt{--paper}, \texttt{--ink}/\texttt{forest},
\texttt{--accent}, \texttt{--muted}, \texttt{--card}, \texttt{--line}).

\subsection{\texttt{eslint.config.mjs}}
\texttt{eslint-config-next}: App Router-regels + TypeScript.

\subsection{\texttt{proxy.ts}}
Next~16 proxy/middleware-equivalent. Matcher: \texttt{/admin} en \texttt{/admin/:path*}.
\begin{itemize}[leftmargin=*]
  \item Niet ingelogd op adminpagina $\to$ redirect \texttt{/admin/login?next=...}
  \item Wel ingelogd op loginpagina $\to$ redirect \texttt{/admin}
  \item \textbf{API-routes} (\texttt{/api/*}) vallen \emph{niet} onder deze matcher;
        elke admin-API herhaalt cookiecheck zelf.
\end{itemize}

\subsection{\texttt{AGENTS.md} / \texttt{README.md}}
Operationele notities voor agents/developers: defaults, geen verplichte \texttt{.env},
waarschuwing om \texttt{node\_modules/}, \texttt{data/*}, \texttt{.next/} niet per ongeluk te committen
(er is geen \texttt{.gitignore} in deze repo).

\section{Map \texttt{lib/} --- domeinlaag}

\subsection{\texttt{lib/config.ts} --- omgeving en seeddefinities}
\begin{itemize}[leftmargin=*]
  \item \texttt{getEvents()} leest seed-events uit (prioriteit):
    \texttt{EVENTS\_JSON} $\to$ legacy \texttt{EVENT\_*} vars $\to$ hardcoded fallback
    (\texttt{afc-avond}, \texttt{afc-workshop}).
  \item \texttt{getAppUrl()} uit \texttt{APP\_URL} (default \texttt{http://localhost:3000}).
  \item \texttt{getAdminPassword()} uit \texttt{ADMIN\_PASSWORD} (default \texttt{admin123}).
  \item \textbf{Belangrijk:} dit is \emph{seed/config}, niet de live eventbron.
        Runtime-UI leest events uit SQLite via \texttt{lib/db.ts}.
\end{itemize}

\subsection{\texttt{lib/session.ts} --- sessietoken}
\begin{lstlisting}
SESSION_COOKIE = "gate_session"
token = hex(SHA-256(UTF8("gate:" + adminPassword)))
\end{lstlisting}
Vergelijking gebeurt length-check + XOR-loop (timing-vriendelijker dan \texttt{===}).
Er zit \textbf{geen} expiry in het token zelf; cookie \texttt{maxAge} (typisch 12u) bepaalt levensduur.
\textbf{Keuze:} simpel, geen JWT-library, werkt in Edge/proxy via Web Crypto.

\subsection{\texttt{lib/auth.ts}}
Re-export van session helpers + \texttt{requireAdmin()} voor Server Components:
ongeldige cookie $\to$ \texttt{redirect("/admin/login")}.

\subsection{\texttt{lib/ticket-token.ts} --- ondertekende tickets}
Formaat: \texttt{t1.<base64url(JSON claims)>.<base64url(HMAC-SHA256(body))>}.

Claims (\texttt{v: 1}): \texttt{id}, \texttt{eventId}, \texttt{eventSlug}, \texttt{eventName},
\texttt{name}, \texttt{email}, \texttt{phone}, \texttt{extraInfo}, \texttt{foodPreference},
\texttt{cvOriginalName}, \texttt{createdAt}.

Secret: \texttt{TICKET\_SECRET} $\to$ anders \texttt{ADMIN\_PASSWORD} $\to$ anders
\texttt{demo-ticket-secret}.

\texttt{verifyTicketToken} gebruikt \texttt{timingSafeEqual}. Bij succes bouwt het een
\texttt{Registration}-achtig object met \texttt{cvStoredName = id + ".pdf"} en
\texttt{checkedInAt = null} (check-in status komt altijd uit DB als die er is).

\textbf{Waarom signed tokens?} Op Vercel is DB-state per instance niet gedeeld. De QR moet
toch verifieerbaar blijven; daarna kan \texttt{getRegistrationByToken} de claims in de lokale DB upserten.

\subsection{\texttt{lib/ticket.ts}}
\begin{itemize}[leftmargin=*]
  \item \texttt{ticketUrl(slug, token)} $\to$ \texttt{\{APP\_URL\}/ticket/\{slug\}/\{encodeURIComponent(token)\}}.
  \item \texttt{parseTicketPayload(raw)} accepteert:
    prefixes \texttt{TICKET:}/\texttt{AFC:}, volledige ticket-URL's (\texttt{/ticket/} of \texttt{/t/}),
    raw signed \texttt{t1...}, of legacy tokens \texttt{[A-Za-z0-9\_-]\{16,\}}.
\end{itemize}
Dit maakt de scanner tolerant voor ``camera leest URL'' versus ``camera leest kale token''.

\subsection{\texttt{lib/qr.ts}}
Genereert QR die de \emph{ticket-URL} encodeert (niet enkel de token), 360px, ECL M,
kleuren navy/wit. \texttt{qrDataUrl} voor UI; \texttt{qrPngBuffer} voor e-mail-CID.

\subsection{\texttt{lib/validation.ts}}
Zod-schema \texttt{registrationFields}: slug, naam (2--120), e-mail (lower+regex),
telefoon (8--40), extraInfo (max 2000, default leeg), foodPreference $\in$ \texttt{FOOD\_OPTIONS}.

\subsection{\texttt{lib/food.ts}}
Enum/lijst voedselvoorkeuren + \texttt{foodLabel()}. Één bron voor formulier, validatie en CSV.

\subsection{\texttt{lib/types.ts}}
\texttt{Registration} (camelCase) $\leftrightarrow$ \texttt{RegistrationRow} (SQL join-kolommen)
via \texttt{mapRegistration}. Voorkomt dat UI per ongeluk snake\_case verwacht.

\subsection{\texttt{lib/datetime.ts}}
\begin{itemize}[leftmargin=*]
  \item Weergave in \texttt{Europe/Brussels} met \texttt{nl-BE}.
  \item \texttt{brusselsLocalToIso}: zet \texttt{datetime-local} (Brussel-wandtijd) om naar UTC ISO
        via iteratieve offsetcorrectie (zonder extra TZ-library).
  \item \texttt{isoToBrusselsInput}: omgekeerd voor adminformulieren.
\end{itemize}

\subsection{\texttt{lib/csv.ts}}
Excel-vriendelijke CSV: UTF-8 BOM, delimiter \texttt{;}, alle cellen gequoted.
Kolommen: Event, Naam, E-mail, Telefoon, Voedselvoorkeur, Extra info, CV,
Ingeschreven op, Aanwezig, Check-in tijd, Ticketcode.

\subsection{\texttt{lib/email.ts}}
Als \texttt{SMTP\_HOST} én \texttt{SMTP\_FROM} ontbreken: mail wordt overgeslagen
(\texttt{\{sent:false, reason:"not-configured"\}}). Anders nodemailer + inline QR (CID).
Inschrijving faalt \textbf{niet} als mail faalt (best-effort).

\subsection{\texttt{lib/db.ts} --- hart van de datalaag}
Zie hoofdstuk~\ref{ch:data}. Bevat schema, migratie single$\to$multi-event, seed,
CRUD voor events/registrations, check-in, file cleanup, Vercel padlogica, en
signed-token upsert bij tokenlookup.

\section{Map \texttt{app/} --- routes en UI}

\subsection{Layout \& globale stijl}
\begin{itemize}[leftmargin=*]
  \item \texttt{app/layout.tsx}: \texttt{lang="nl"}, Josefin Sans CSS-variable, metadata ``AFC Ticketing''.
  \item \texttt{app/globals.css}: design tokens + utility helpers + printregels.
  \item \texttt{app/not-found.tsx}: generieke 404 (o.a.\ onbekend ticket/event).
\end{itemize}

\subsection{Publieke pages}
\begin{tabularx}{\textwidth}{@{}lX@{}}
\toprule
\textbf{Pad} & \textbf{Bestand / rol} \\
\midrule
\texttt{/} & \texttt{app/page.tsx} --- open events (\texttt{getOpenEvents}), optionele foto \\
\texttt{/events/[slug]} & inschrijfpagina: intro + optionele \texttt{description} + \texttt{RegisterForm} \\
\texttt{/ticket/[eventSlug]/[token]} & ticket + QR; query \texttt{nieuw}, \texttt{mail} \\
\texttt{/afgelopen} & gearchiveerde events (read-only) \\
\bottomrule
\end{tabularx}

Alle list/detail pages gebruiken \texttt{export const dynamic = "force-dynamic"} zodat SQLite-wijzigingen
meteen zichtbaar zijn (geen statische cache van eventlijsten).

\subsection{Admin pages}
\begin{tabularx}{\textwidth}{@{}lX@{}}
\toprule
\textbf{Pad} & \textbf{Rol} \\
\midrule
\texttt{/admin/login} & wachtwoordformulier (\texttt{LoginForm}) \\
\texttt{/admin} & dashboard stats + snelle links \\
\texttt{/admin/events} & \texttt{EventManager} CRUD \\
\texttt{/admin/deelnemers} & \texttt{AttendeeTable} (+ \texttt{?event=}) \\
\texttt{/admin/scanner} & \texttt{CheckinScanner} \\
\bottomrule
\end{tabularx}
Elke admin page (behalve login) start met \texttt{await requireAdmin()}.

\section{Map \texttt{components/}}
\begin{description}[style=nextline,leftmargin=0pt]
  \item[\texttt{SiteHeader.tsx}] Publieke header met AFC-logo + link naar organisatie/login.
  \item[\texttt{AdminNav.tsx}] Clientnav + logout (\texttt{DELETE /api/auth}).
  \item[\texttt{LoginForm.tsx}] \texttt{POST /api/auth} met password; volgt \texttt{?next=}.
  \item[\texttt{RegisterForm.tsx}] Client form $\to$ \texttt{POST /api/register} (multipart) $\to$ redirect ticket.
  \item[\texttt{TicketCard.tsx}] Toont gegevens + QR + print; waarschuwt als mail faalde.
  \item[\texttt{EventManager.tsx}] Admin eventform: intro, uitgebreide beschrijving, optionele foto,
        open/close venster, archiveren, verwijderen.
  \item[\texttt{AttendeeTable.tsx}] Zoeken, filteren, manueel toevoegen/verwijderen, check-in toggle, CV-link, export.
  \item[\texttt{CheckinScanner.tsx}] Camera + manuele token/e-mail check-in via \texttt{/api/checkin}.
\end{description}

\section{Map \texttt{public/}}
Relevant: \texttt{afc-gent-logo.svg}. Overige SVG's zijn Next-starterresten en niet functioneel nodig.

\section{Map \texttt{data/} (runtime)}
\begin{lstlisting}
data/
  tickets.db          # SQLite hoofdbestand
  tickets.db-wal/shm  # WAL sidecars (lokaal)
  uploads/
    {registrationId}.pdf
    event-{eventId}.jpg|png|webp|gif
\end{lstlisting}
Op Vercel: zelfde structuur onder \texttt{/tmp/ticketing-data/}.

%======================================================================
\chapter{Datamodel en persistentie}
\label{ch:data}
%======================================================================

\section{Initialisatievolgorde bij cold start}
\begin{enumerate}[leftmargin=*]
  \item Bepaal data-dir (cwd/\texttt{data} of \texttt{/tmp/ticketing-data}).
  \item Maak uploadmap aan.
  \item Open SQLite; journal mode WAL (lokaal) of DELETE (Vercel).
  \item \texttt{foreign\_keys = ON}.
  \item \texttt{ensureEventSchema} (\texttt{CREATE TABLE IF NOT EXISTS} + \texttt{ensureColumn}).
  \item Seed events als tabel leeg is (\texttt{getEvents()} uit config).
  \item Migreer legacy single-event \texttt{registrations} indien nodig.
  \item Zorg dat registrations-tabel + indexes bestaan.
  \item Cache connectie op \texttt{globalThis} voor hot reload / warm serverless.
\end{enumerate}

\section{Tabel \texttt{events}}
\begin{lstlisting}
id TEXT PRIMARY KEY
slug TEXT NOT NULL UNIQUE
name TEXT NOT NULL
date TEXT NOT NULL                 -- ISO timestamp
location TEXT NOT NULL
intro TEXT NOT NULL                -- korte tekst homepage
description TEXT NOT NULL DEFAULT '' -- enkel eventpagina
image_stored_name TEXT              -- optioneel bestand in uploads/
is_open INTEGER NOT NULL DEFAULT 1  -- handmatige schakelaar
archived INTEGER NOT NULL DEFAULT 0
registration_opens_at TEXT          -- optioneel ISO
registration_closes_at TEXT         -- optioneel ISO
created_at TEXT NOT NULL
\end{lstlisting}

\subsection{Acceptatieregel inschrijving}
Een event accepteert inschrijvingen iff:
\[
\neg archived \;\wedge\; is\_open=1 \;\wedge\;
(now \ge opens\_at \lor opens\_at=\emptyset) \;\wedge\;
(now \le closes\_at \lor closes\_at=\emptyset).
\]
Homepage toont enkel events die deze regel passeren (\texttt{getOpenEvents}).

\section{Tabel \texttt{registrations}}
\begin{lstlisting}
id TEXT PRIMARY KEY
event_id TEXT NOT NULL REFERENCES events(id) ON DELETE CASCADE
name TEXT NOT NULL
email TEXT NOT NULL COLLATE NOCASE
phone TEXT NOT NULL
extra_info TEXT NOT NULL DEFAULT ''
food_preference TEXT NOT NULL
cv_original_name TEXT NOT NULL
cv_stored_name TEXT NOT NULL
ticket_token TEXT NOT NULL UNIQUE
checked_in_at TEXT                  -- NULL = nog niet binnen
created_at TEXT NOT NULL
UNIQUE(event_id, email)
\end{lstlisting}
Indexes op \texttt{event\_id}, \texttt{checked\_in\_at}, \texttt{created\_at}.

\section{Migratie single-event $\to$ multi-event}
Oude DB's hadden \texttt{registrations} zonder \texttt{event\_id}.
Procedure:
\begin{enumerate}[leftmargin=*]
  \item Detecteer via \texttt{PRAGMA table\_info}.
  \item Rename naar \texttt{registrations\_legacy} (met herstelpad als een vorige run crashte en de legacy-tabel bleef bestaan).
  \item Maak nieuwe tabel; \texttt{INSERT OR IGNORE ... SELECT} met fallback \texttt{event\_id} = oudste event.
  \item Drop legacy.
\end{enumerate}
Deze migratie is bewust idempotent/resilient: een achtergebleven
\texttt{registrations\_legacy} mag \texttt{npm run build}/start niet breken.

\section{Bestanden op schijf}
\begin{itemize}[leftmargin=*]
  \item CV: \texttt{\{registrationId\}.pdf}
  \item Eventfoto: \texttt{event-\{eventId\}.\{ext\}}
  \item Bij image replace/delete of event/registration delete: best-effort \texttt{unlink}.
\end{itemize}

\section{Domain mapping}
SQL join levert \texttt{event\_slug}/\texttt{event\_name}; \texttt{mapRegistration} zet om naar camelCase.
UI en API werken bij voorkeur met het domain model; DB-helpers accepteren vaak beide lagen intern.

%======================================================================
\chapter{Authenticatie en autorisatie}
%======================================================================

\section{Admin login flow}
\begin{enumerate}[leftmargin=*]
  \item UI: \texttt{POST /api/auth} JSON \texttt{\{password\}}.
  \item Server vergelijkt met \texttt{getAdminPassword()}.
  \item Bij succes: cookie \texttt{gate\_session=SHA256("gate:"+password)},
        \texttt{httpOnly}, \texttt{sameSite=lax}, \texttt{path=/},
        \texttt{secure} als \texttt{APP\_URL} met \texttt{https} begint, \texttt{maxAge} $\approx$ 12u.
  \item Logout: \texttt{DELETE /api/auth} verwijdert cookie.
\end{enumerate}

\section{Beschermingslagen (defense in depth)}
\begin{enumerate}[leftmargin=*]
  \item \texttt{proxy.ts} blokkeert admin \emph{pages}.
  \item \texttt{requireAdmin()} in server pages (extra redirect).
  \item Elke gevoelige API checkt opnieuw de cookie (want \texttt{/api} zit niet in proxy matcher).
\end{enumerate}
Publiek blijven o.a.\ \texttt{/api/register} en \texttt{/api/event-image/[id]}.
CV-download is admin-only.

\section{Beperkingen van dit authmodel}
\begin{itemize}[leftmargin=*]
  \item Geen per-user accounts, geen audit log van wie inlogde.
  \item Token verandert automatisch als adminwachtwoord verandert (alle sessies invalid).
  \item Geen CSRF-token; mitigatie via \texttt{SameSite=Lax} + same-origin form posts.
  \item Voor internet-facing productie: sterk wachtwoord via env, HTTPS, reverse proxy.
\end{itemize}

%======================================================================
\chapter{API-referentie}
%======================================================================

Alle hieronder genoemde routes draaien in \texttt{runtime = "nodejs"} (nodig voor
\texttt{better-sqlite3} en filesystem).

\section{Publiek}

\subsection{\texttt{POST /api/register}}
Multipart form:
\texttt{eventSlug}, \texttt{name}, \texttt{email}, \texttt{phone}, \texttt{foodPreference},
\texttt{extraInfo?}, \texttt{cv} (PDF, verplicht, $\le$ 5\,MB).

Stappen:
\begin{enumerate}[leftmargin=*]
  \item Zod-parse $\to$ 400 bij fout.
  \item Event bestaat + accepteert inschrijvingen $\to$ anders 404.
  \item PDF-type/grootte $\to$ 400/413.
  \item Unieke e-mail per event $\to$ anders 409.
  \item \texttt{randomUUID()} als id; schrijf CV; mint HMAC-token; insert DB.
  \item Best-effort mail.
  \item JSON \texttt{\{ok, eventSlug, token, emailSent\}} (UI redirect naar ticket).
\end{enumerate}

\subsection{\texttt{GET /api/event-image/[id]}}
Publiek. 404 als geen image. Stream bestand met content-type op extensie; cachebaar.

\section{Auth}
\subsection{\texttt{POST /api/auth}} / \texttt{DELETE /api/auth}
Zie vorig hoofdstuk.

\section{Admin --- check-in \& export}

\subsection{\texttt{POST /api/checkin}}
Cookie verplicht. Body JSON kan bevatten:
\texttt{payload} (scanruwe tekst), \texttt{token}, \texttt{id}, \texttt{email},
\texttt{eventSlug?}, \texttt{undo?}.

Gedrag:
\begin{itemize}[leftmargin=*]
  \item \texttt{undo:true} + id: clear \texttt{checked\_in\_at}.
  \item Anders resolve registratie; zet check-in timestamp als nog leeg.
  \item Ambiguous e-mail over events $\to$ 409 met matches.
  \item Onbekend $\to$ 404.
  \item Succes: naam, foodPreference, alreadyCheckedIn, eventName, \ldots
\end{itemize}

\subsection{\texttt{GET /api/export?event=slug?}}
CSV download; optioneel gefilterd op event.

\subsection{\texttt{GET /api/cv/[id]}}
Admin-only PDF stream met originele filename in \texttt{Content-Disposition}.

\section{Admin --- events \texttt{/api/admin/events}}
\begin{tabularx}{\textwidth}{@{}lX@{}}
\toprule
\textbf{Method} & \textbf{Body / effect} \\
\midrule
POST & multipart create (+ optionele image) \\
PUT & multipart update (id + velden + optionele image/clearImage) \\
PATCH & JSON \texttt{\{id, archived\}} \\
DELETE & JSON \texttt{\{id\}} cascade registrations + files \\
\bottomrule
\end{tabularx}
Velden o.a.\ name, slug?, date (datetime-local Brussel), location, intro, description,
isOpen, archived, registrationOpensAt/ClosesAt, image, clearImage.
Images: jpeg/png/webp/gif, max 5\,MB.

\section{Admin --- attendees \texttt{/api/admin/attendees}}
\begin{itemize}[leftmargin=*]
  \item \textbf{POST} multipart: zelfde velden als register; CV optioneel;
        mag ook op gesloten events; optioneel \texttt{sendEmail=1}.
  \item \textbf{DELETE} JSON \texttt{\{id\}} + CV cleanup.
\end{itemize}

%======================================================================
\chapter{End-to-end flows}
%======================================================================

\section{Publieke inschrijving}
\begin{lstlisting}
Browser GET /
  -> Server: getOpenEvents() from SQLite
  -> Cards (image? + intro)
User clicks event
  -> GET /events/{slug}
  -> intro + description? + RegisterForm
User submits multipart
  -> POST /api/register
  -> validate, store PDF, mint t1-token, insert row, maybe SMTP
  -> redirect /ticket/{slug}/{token}?nieuw=1&mail=0|1
Ticket page
  -> getRegistrationByToken (DB or verify+upsert)
  -> qrDataUrl(ticketUrl)
  -> TicketCard
\end{lstlisting}

\section{Deurcheck-in}
\begin{lstlisting}
Admin login -> /admin/scanner
Camera reads QR (URL or token)
  -> POST /api/checkin { payload }
  -> parseTicketPayload -> getRegistrationByToken -> set checked_in_at
UI shows name + food preference + alreadyCheckedIn?
\end{lstlisting}

\section{Admin event met foto + lange beschrijving}
\begin{lstlisting}
/admin/events -> EventManager
Create/Edit form includes:
  intro (homepage), description (event page only),
  optional image upload / clearImage
POST/PUT /api/admin/events
  -> create/update row; save event-{id}.ext
Homepage shows image if image_stored_name set
Event page shows description if non-empty
\end{lstlisting}

\section{Export}
Admin klikt export $\to$ \texttt{GET /api/export} $\to$ CSV met BOM voor Excel.

%======================================================================
\chapter{UI-architectuur en designkeuzes}
%======================================================================

\section{Server vs client}
Server Components lezen DB direct (geen waterfalls via client fetch voor initial paint).
Client Components isoleren browser-API's (File input, camera, \texttt{useRouter.refresh()}).

\section{Visuele taal}
Navy \texttt{\#081C3C}, accent oranje \texttt{\#EC6525}, papierachtige achtergrond,
grote serif/display headings via utility \texttt{.serif}, afgeronde kaarten ($\approx$ 28px).
Doel: herkenbaar AFC/organisatiegevoel zonder zwaar design system.

\section{Waarom geen aparte componentlibrary}
Voor deze scope (enkele tientallen schermen) zijn Tailwind + lokale componenten sneller
dan shadcn/MUI te onderhouden. Interactieve ``cards'' bestaan alleen als echte UI-containers
(formulieren, eventkaarten), niet als decoratie.

%======================================================================
\chapter{Omgeving, deploy en operationele concerns}
%======================================================================

\section{Environment variables}
\begin{tabularx}{\textwidth}{@{}lX@{}}
\toprule
\textbf{Var} & \textbf{Rol / default} \\
\midrule
\texttt{ADMIN\_PASSWORD} & admin login (default \texttt{admin123}) \\
\texttt{APP\_URL} & absolute URL in QR/mail (default localhost:3000) \\
\texttt{TICKET\_SECRET} & HMAC secret (valt terug op admin password) \\
\texttt{EVENTS\_JSON} / \texttt{EVENT\_*} & seed bij lege DB \\
\texttt{SMTP\_HOST}, \texttt{SMTP\_FROM}, \ldots & optionele mail \\
\texttt{VERCEL}/\texttt{VERCEL\_ENV} & schakelt data-dir naar \texttt{/tmp} \\
\bottomrule
\end{tabularx}
Geen verplichte \texttt{.env} voor lokale demo.

\section{Scripts}
\begin{lstlisting}
npm install
npm run dev      # http://localhost:3000
npm run lint
npm run build && npm start
\end{lstlisting}

\section{VPS (aanbevolen productie)}
Persistente schijf, Node 22, reverse proxy met HTTPS, env met sterk wachtwoord +
\texttt{APP\_URL} + SMTP. Backup: stop service, kopieer \texttt{data/}.
Update: \texttt{git pull \&\& npm ci \&\& npm run build \&\& systemctl restart \ldots}.

\section{Vercel-beperkingen}
Geschikt als preview/demo. Schrijfacties in \texttt{/tmp} verdwijnen tussen cold starts.
Signed tickets mitigeren ticketweergave; blijvende registraties vereisen persistente storage
(dus VPS of externe object/DB store --- niet in huidige code).

\section{Git-hygiëne}
Repo mist \texttt{.gitignore}. Commit nooit per ongeluk \texttt{node\_modules/},
\texttt{.next/}, \texttt{data/tickets.db*}, of grote uploadbinaries.

%======================================================================
\chapter{Reconstructiegids: bouw de tool opnieuw}
%======================================================================

\section{Fase 0 --- scaffold}
\begin{enumerate}[leftmargin=*]
  \item Next.js App Router + TS + Tailwind 4 + ESLint.
  \item Dependencies: better-sqlite3, zod, qrcode, nodemailer, react-qr-scanner + types.
  \item Path alias \texttt{@/*}; \texttt{proxy.ts} voor \texttt{/admin/*}.
\end{enumerate}

\section{Fase 1 --- datalaag}
Implementeer \texttt{lib/db.ts} exact volgens schema hierboven, inclusief seed + migratie.
Schrijf eerst helpers: \texttt{createEvent}, \texttt{createRegistration},
\texttt{getRegistrationByToken} (met verify+upsert), \texttt{checkInRegistration}.

\section{Fase 2 --- tokens \& validatie}
\texttt{ticket-token.ts}, \texttt{ticket.ts}, \texttt{validation.ts}, \texttt{food.ts}, \texttt{datetime.ts}.

\section{Fase 3 --- publieke API + UI}
\texttt{POST /api/register}, homepage, eventpagina, ticketpagina, optionele mail/QR.

\section{Fase 4 --- admin}
Auth cookie, dashboard, events CRUD (+ image/description), attendees, scanner, export, CV download.

\section{Fase 5 --- harding}
force-dynamic op data pages; Vercel padlogica; image/CV size limits; unique email constraint;
archiefvensters; resilient legacy migratie.

\section{Acceptatiechecklist}
\begin{itemize}[leftmargin=*]
  \item Inschrijven mét PDF werkt; duplicate e-mail $\to$ 409.
  \item Ticket QR opent dezelfde ticketpagina; scanner checkt in.
  \item Event zonder foto/beschrijving blijft netjes; mét foto/beschrijving toont correcte plaatsen.
  \item Zonder SMTP slaagt register toch.
  \item Admin delete event ruimt registrations + files op.
  \item \texttt{npm run lint} en \texttt{npm run build} groen.
\end{itemize}

%======================================================================
\chapter{Beslissingen, trade-offs en uitbreidingspaden}
%======================================================================

\section{Belangrijkste trade-offs}
\begin{itemize}[leftmargin=*]
  \item \textbf{SQLite vs Postgres:} eenvoud vs horizontale schaal.
  \item \textbf{Signed tickets vs opaque DB tokens:} serverless-resilience vs revoke-from-DB-only.
    Huidige hybride: opaque unique kolom + HMAC claims.
  \item \textbf{Shared admin password vs echte users:} snelle ops vs accounting.
  \item \textbf{CV op lokale schijf vs S3:} eenvoud vs multi-instance deeling.
  \item \textbf{Geen testsuite:} snellere iteratie; regressierisico --- overweeg later API integration tests.
\end{itemize}

\section{Mogelijke evoluties (niet geïmplementeerd)}
\begin{itemize}[leftmargin=*]
  \item Object storage voor CV's/images.
  \item Postgres + migrations framework.
  \item Rate limiting op \texttt{/api/register}.
  \item Per-scanner accounts / role-based access.
  \item Soft-delete en audit trail.
  \item Officiële \texttt{.gitignore} en \texttt{.env.example}.
\end{itemize}

%======================================================================
\chapter{Appendix}
%======================================================================

\section{Snelle bestand-naar-verantwoordelijkheid mapping}
\begin{longtable}{@{}p{0.42\textwidth}p{0.5\textwidth}@{}}
\toprule
Bestand & Verantwoordelijkheid \\
\midrule
\endhead
\texttt{lib/db.ts} & Schema, CRUD, migratie, uploads pad \\
\texttt{lib/config.ts} & Env + seed events \\
\texttt{lib/session.ts}/\texttt{auth.ts} & Cookie auth \\
\texttt{lib/ticket-token.ts} & HMAC tickets \\
\texttt{lib/ticket.ts}/\texttt{qr.ts} & URL + QR \\
\texttt{lib/validation.ts}/\texttt{food.ts} & Inputcontracten \\
\texttt{lib/email.ts}/\texttt{csv.ts}/\texttt{datetime.ts} & Mail, export, tijd \\
\texttt{proxy.ts} & Admin page gate \\
\texttt{app/api/register/route.ts} & Publieke inschrijving \\
\texttt{app/api/auth/route.ts} & Login/logout \\
\texttt{app/api/checkin/route.ts} & Check-in/undo \\
\texttt{app/api/export/route.ts} & CSV \\
\texttt{app/api/cv/[id]/route.ts} & CV download \\
\texttt{app/api/event-image/[id]/route.ts} & Publieke eventfoto \\
\texttt{app/api/admin/events/route.ts} & Eventbeheer \\
\texttt{app/api/admin/attendees/route.ts} & Deelnemers add/delete \\
\texttt{components/EventManager.tsx} & Admin event UI \\
\texttt{components/AttendeeTable.tsx} & Admin deelnemers UI \\
\texttt{components/CheckinScanner.tsx} & Scanner UI \\
\texttt{components/RegisterForm.tsx} & Publiek formulier \\
\texttt{components/TicketCard.tsx} & Ticketweergave \\
\bottomrule
\end{longtable}

\section{HTTP-statuscodes die de API consequent gebruikt}
\begin{itemize}[leftmargin=*]
  \item 200 --- succes
  \item 400 --- validatiefout / slechte input
  \item 401 --- niet ingelogd
  \item 404 --- onbekend event/ticket/deelnemer / gesloten inschrijving
  \item 409 --- conflict (duplicate e-mail, ambiguous check-in)
  \item 413 --- upload te groot
\end{itemize}

\section{Slot}
Wie dit document volgt, kan de AFC Ticketing-tool opnieuw opzetten met dezelfde
architectuurkeuzes: monolithische Next.js-app, SQLite, HMAC-tickets, optionele SMTP,
multi-event admin, en een bewuste scheiding tussen korte intro (homepage) en uitgebreide
beschrijving + optionele foto (eventkaart/eventpagina).

\end{document}
